\documentclass[11pt,a4paper]{article}
\usepackage[T1]{fontenc}
\usepackage[utf8]{inputenc}
\usepackage[brazil]{babel}
\usepackage{amsmath,amssymb}
\usepackage[margin=2.6cm]{geometry}
\usepackage{graphicx}
\usepackage{xcolor}
\usepackage[colorlinks=true,linkcolor=blue!50!black,citecolor=blue!50!black,
            urlcolor=blue!50!black]{hyperref}
\usepackage{mdframed}
\usepackage{enumitem}

\newcommand{\Msun}{M_\odot}
\newcommand{\Mch}{M_{\rm Ch}}
\newcommand{\code}[1]{\texttt{\small #1}}

\newmdenv[linecolor=blue!40,backgroundcolor=blue!3,linewidth=0.6pt,
          skipabove=8pt,skipbelow=8pt]{caixa}
\newmdenv[linecolor=orange!60,backgroundcolor=orange!4,linewidth=0.6pt,
          skipabove=8pt,skipbelow=8pt]{alerta}

\newcounter{exerc}[section]
\newcommand{\ex}[1]{\refstepcounter{exerc}%
  \medskip\noindent\textbf{Exercício \thesection.\theexerc.}\ #1\medskip}

\title{\bf Anãs brancas magnetizadas super-Chandrasekhar\\
\large Uma apostila para começar}
\author{Grupo WD-MAG --- CCT/UDESC Joinville}
\date{\today}

\begin{document}
\maketitle

\begin{abstract}
\noindent
Esta apostila reúne o que é preciso saber para começar a trabalhar no projeto
\code{wd-magnetizada}. Ela parte de matéria degenerada e chega até as
simulações 3D de magnetoidrodinâmica que o grupo roda em cluster, passando pela
construção de equilíbrios e pela teoria de estabilidade de campos estelares.
Não pressupõe nada além de mecânica, eletromagnetismo e termodinâmica de
graduação.

O texto é \emph{didático e incompleto de propósito}: cada seção termina
apontando para onde está o tratamento rigoroso, seja na literatura, seja no
arquivo \code{docs/teoria.md} do repositório, seja no próprio código. Os
exercícios não são decorativos --- vários deles reproduzem números que o grupo
usa no dia a dia, e resolvê-los é a forma mais rápida de saber se você entendeu.
\end{abstract}

\tableofcontents
\newpage

%============================================================
\part*{Parte I --- De onde vem o problema}
\addcontentsline{toc}{part}{Parte I --- De onde vem o problema}

\section{Matéria degenerada e o limite de Chandrasekhar}

\subsection{Por que uma anã branca não colapsa}

Uma anã branca é o que sobra de uma estrela de massa baixa ou intermediária
depois que a fusão termina. Sem fusão não há pressão térmica sustentada, e a
estrela deveria colapsar sob a própria gravidade. O que a segura é um efeito
puramente quântico: o \textbf{princípio de exclusão de Pauli}.

Elétrons são férmions. Comprima o gás e eles não podem ocupar o mesmo estado;
são forçados a momentos cada vez maiores, e esse momento gera pressão mesmo a
temperatura zero. É a \emph{pressão de degenerescência}.

Num gás completamente degenerado a $T=0$, todos os estados abaixo do momento de
Fermi $p_F$ estão ocupados e nenhum acima. Contando estados no espaço de fase,
\begin{equation}
n_e = \frac{8\pi}{3h^3}\,p_F^3
\qquad\Longrightarrow\qquad
p_F = \hbar\,(3\pi^2 n_e)^{1/3}.
\end{equation}

\begin{caixa}
\textbf{O parâmetro que organiza tudo.} Define-se
\begin{equation}
x \equiv \frac{p_F}{m_e c},
\end{equation}
a razão entre o momento de Fermi e $m_e c$. É o parâmetro de
relativisticidade: $x \ll 1$ é gás não relativístico, $x \gg 1$ é
ultrarrelativístico. Toda a equação de estado se escreve em termos dele.
\end{caixa}

\subsection{A equação de estado de Chandrasekhar}

Integrando a pressão do gás de Fermi sobre a esfera de Fermi, com a relação
relativística $E^2 = (pc)^2 + (m_ec^2)^2$, obtém-se a forma fechada
\begin{equation}
\boxed{\;
P(x) = A\left[x(2x^2-3)\sqrt{x^2+1} + 3\,\mathrm{arcsinh}\,x\right],
\qquad
\rho(x) = B\,x^3 \;}
\label{eq:eos}
\end{equation}
com
\begin{equation}
A = \frac{\pi m_e^4 c^5}{3h^3} = 6{,}01\times10^{22}\ \mathrm{dyn\,cm^{-2}},
\qquad
B = 9{,}82\times10^5\,\mu_e\ \mathrm{g\,cm^{-3}} .
\end{equation}

Aqui $\mu_e$ é o número de núcleons por elétron. Para matéria de carbono-oxigênio
completamente ionizada, $\mu_e = 2$: cada elétron ``carrega'' dois núcleons. É o
valor usado em todo o projeto.

\begin{caixa}
\textbf{Onde isso está no código.} \code{scf/eos.py}: a função
\code{pressure(x)} é a Eq.~\eqref{eq:eos}, \code{density(x, mu\_e)} é
$\rho(x)$, \code{B\_of\_mu\_e(mu\_e)} devolve $B$. Abra o arquivo --- são
cem linhas e você reconhece cada fórmula.
\end{caixa}

\subsection{A entalpia, e por que ela importa}

Para construir equilíbrios vamos precisar da \emph{entalpia específica}
\begin{equation}
H \equiv \int \frac{dP}{\rho} .
\end{equation}
Para a EOS acima a integral sai analiticamente, o que é a razão técnica de todo
o método numérico do Capítulo~4 funcionar:
\begin{equation}
\boxed{\;H(x) = \frac{8A}{B}\left[\sqrt{1+x^2} - 1\right]\;}
\label{eq:entalpia}
\end{equation}
e a inversa, que é o que de fato se usa,
\begin{equation}
x(H) = \sqrt{\left(1 + \frac{HB}{8A}\right)^2 - 1},
\qquad H<0 \Rightarrow \rho = 0 .
\label{eq:inversa}
\end{equation}
O $-1$ normaliza $H=0$ na superfície, onde $\rho \to 0$. Isso define a
superfície da estrela sem precisar procurá-la: é o conjunto $H = 0$.

\ex{Verifique que \eqref{eq:inversa} é realmente a inversa de
\eqref{eq:entalpia}. Depois mostre que para $x \ll 1$ vale
$\rho \propto H^{3/2}$ e para $x \gg 1$ vale $\rho \propto H^{3}$.}

\subsection{O limite de Chandrasekhar}

O exercício anterior é a chave. Uma relação $\rho \propto H^{n}$ corresponde a
um politropo de índice $n$. A EOS degenerada \emph{desliza} entre $n=3/2$ no
regime não relativístico e $n=3$ no ultrarrelativístico.

E $n=3$ é o caso marginal: para um politropo de índice $3$, a massa total
\textbf{não depende} da densidade central. Comprimir mais não adiciona massa.
Esse é o limite de Chandrasekhar,
\begin{equation}
\boxed{\;\Mch = \frac{5{,}816}{\mu_e^2}\,\Msun \simeq 1{,}44\,\Msun
\quad (\mu_e = 2).\;}
\end{equation}

\begin{alerta}
\textbf{Leitura correta do limite.} $\Mch$ não é uma proibição mágica. É o
valor assintótico de uma sequência de equilíbrios construída com uma física
específica: gás de elétrons degenerado, sem rotação, sem campo, sem
temperatura. Mude a física e o número muda. \emph{Todo o projeto é sobre
quanto ele muda e a que custo.}
\end{alerta}

\ex{A $\rho_c = 10^9$ g/cm$^3$ e $\mu_e=2$, uma anã branca sem rotação e sem
campo pesa $1{,}3464\,\Msun$. A $5\times10^9$ pesa $1{,}4012$. Estime, ajustando
uma aproximação simples, a que $\rho_c$ ela chegaria a $1{,}44$. Compare com o
limiar de neutronização, $1{,}94\times10^{10}$ g/cm$^3$. O que você conclui?}

\section{Por que alguém quer passar do limite}

\subsection{A motivação observacional}

Algumas supernovas do tipo Ia são \emph{superluminosas} demais para o cenário
padrão. Se a luminosidade vem do decaimento de $^{56}$Ni sintetizado na
explosão, e a explosão consome a anã branca inteira, então uma SN Ia
excepcionalmente brilhante sugere uma progenitora excepcionalmente massiva.
O caso clássico é a SN\,2003fg, que pede da ordem de $2\,\Msun$.

Daí a pergunta que organiza toda a literatura que vamos ler: \emph{o que
poderia sustentar uma anã branca acima de $1{,}44\,\Msun$?}

\subsection{As duas rotas, e a contabilidade do virial}

Há duas candidatas clássicas, e o teorema do virial é a forma limpa de
comparar as duas. Para uma configuração estacionária e autogravitante,
\begin{equation}
\boxed{\;2T + W + 3\int P\,dV + \mathcal{E}_{\rm mag} = 0\;}
\label{eq:virial}
\end{equation}
onde $T$ é a energia cinética de rotação, $W < 0$ a energia gravitacional,
e $\mathcal{E}_{\rm mag}$ a energia magnética.

A leitura é direta: rotação e campo entram na mesma equação que a pressão, com
o mesmo sinal, \emph{ambos ajudando a segurar a estrela}. As duas figuras de
mérito que você vai ver o tempo todo no projeto são as frações
\begin{equation}
\frac{T}{|W|}
\qquad\text{e}\qquad
\frac{\mathcal{E}_{\rm mag}}{|W|} .
\end{equation}

\begin{caixa}
\textbf{Ordens de grandeza do nosso modelo de produção.}
$T/|W| = 0{,}099$ e $\mathcal{E}_{\rm mag}/|W| = 0{,}011$. Ou seja: a rotação
carrega dez vezes mais que o campo. Guarde esse par de números --- boa parte
das conclusões do grupo sai dele.
\end{caixa}

\ex{Use \eqref{eq:virial} para argumentar, sem fazer contas detalhadas, por
que uma fração $T/|W|$ da ordem de $0{,}1$ pode alterar a massa
significativamente. Depois estime quanto campo, em gauss, seria preciso numa
estrela de $R \sim 10^8$\,cm e $M \sim 2\,\Msun$ para que
$\mathcal{E}_{\rm mag}/|W| \sim 0{,}1$.}

\subsection{O que limita cada rota}

\textbf{Rotação} é limitada por instabilidades dinâmicas. Quando $T/|W|$ cresce,
a estrela fica suscetível a modos não axissimétricos: o \emph{modo barra} tem
limiar secular em $T/|W| \approx 0{,}14$ e dinâmico em $\approx 0{,}27$. Acima
disso a estrela deforma-se e perde momento angular.

\textbf{Campo} é limitado por duas coisas distintas, e distingui-las é
provavelmente o conceito mais importante desta apostila:
\begin{enumerate}[label=(\alph*)]
\item a estrela precisa \emph{conseguir segurar} o campo --- se a pressão
      magnética supera a do gás em algum ponto, a EOS usada para construir o
      equilíbrio deixa de descrever a própria configuração;
\item a configuração do campo precisa ser \emph{estável} --- e a maioria das
      geometrias simples não é.
\end{enumerate}
O Capítulo~7 trata de (b) e o Capítulo~8 de (a).

%============================================================
\part*{Parte II --- Construindo a estrela}
\addcontentsline{toc}{part}{Parte II --- Construindo a estrela}

\section{Equilíbrio hidrostático e a integral de Bernoulli}

Numa estrela estacionária, a soma das forças por unidade de volume se anula:
\begin{equation}
-\frac{\nabla P}{\rho} - \nabla\Phi
+ \Omega^2(\varpi)\,\varpi\,\hat{e}_\varpi
+ \frac{(\nabla\times\mathbf{B})\times\mathbf{B}}{4\pi\rho} = 0 ,
\label{eq:hidro}
\end{equation}
com $\Phi$ o potencial gravitacional e $\varpi$ a distância ao eixo.

O truque que torna isso tratável: se cada termo for o gradiente de alguma
função escalar, a equação vetorial vira uma equação escalar. O primeiro termo é
$\nabla H$ pela definição de entalpia. O terceiro é $\nabla C_{\rm rot}$ se
$\Omega$ depende só de $\varpi$. O quarto exige mais cuidado (Capítulo~6).
Quando todos são gradientes, integra-se:
\begin{equation}
\boxed{\;H + \Phi - C_{\rm rot} - M_{\rm pol} = C = \text{const.}\;}
\label{eq:bernoulli}
\end{equation}
Essa é a \textbf{integral de Bernoulli}, e é a equação central do método
numérico.

\subsection{O termo de rotação}

Para a lei de rotação diferencial chamada \emph{$j$-constante},
\begin{equation}
\Omega(\varpi) = \frac{\Omega_c A^2}{A^2 + \varpi^2},
\end{equation}
a integral sai em forma fechada:
\begin{equation}
C_{\rm rot}(\varpi) = \int_0^\varpi \Omega^2 \varpi'\,d\varpi'
= \frac{\Omega_c^2 A^2}{2}\,\frac{\varpi^2}{A^2+\varpi^2}.
\end{equation}
O parâmetro $A$ controla o quão diferencial é a rotação: $A \to \infty$
recupera rotação rígida.

\ex{Verifique a integral acima. Depois mostre que $A\to\infty$ dá
$C_{\rm rot}\to \frac{1}{2}\Omega_c^2\varpi^2$, o potencial centrífugo usual.}

\section{O método SCF de Hachisu}

Como resolver \eqref{eq:bernoulli} se $\Phi$ depende de $\rho$, que depende de
$H$, que depende de $\Phi$? Iterando. O método \emph{self-consistent field}
(SCF) de Hachisu (1986) faz:

\begin{enumerate}
\item Chute um $\rho(r,\theta)$ inicial.
\item Resolva a equação de Poisson $\nabla^2\Phi = 4\pi G\rho$ para achar
      $\Phi$.
\item Calcule os termos explícitos ($C_{\rm rot}$, $M_{\rm pol}$).
\item Fixe a constante $C$ impondo a densidade central desejada.
\item Obtenha $H$ de \eqref{eq:bernoulli} e depois $\rho$ de
      \eqref{eq:inversa}.
\item Volte a (2) até $\rho$ parar de mudar.
\end{enumerate}

A equação de Poisson é resolvida por expansão em polinômios de Legendre, que
transforma a integral tridimensional numa soma de integrais radiais:
\begin{equation}
\Phi(r,\theta) = -4\pi G\sum_{l} P_l(\cos\theta)
\left[\frac{1}{r^{l+1}}\!\int_0^r\! \rho_l r'^{\,l+2}dr'
+ r^l\!\int_r^\infty\! \frac{\rho_l}{r'^{\,l-1}}dr'\right].
\end{equation}

\begin{caixa}
\textbf{No código.} \code{scf/poisson.py} implementa exatamente isso, e
\code{scf/scf.py::hachisu\_scf()} é o laço. Vale ler o laço inteiro de uma
sentada --- são umas 60 linhas e é a espinha dorsal de tudo.
\end{caixa}

\ex{Rode o SCF sem rotação e sem campo para $\rho_c = 10^9$ e confira que dá
$1{,}3464\,\Msun$. Depois varie $\rho_c$ e reproduza a curva $M(\rho_c)$ do
Exercício 1.2.}

\section{O campo magnético}

\subsection{A função de fluxo}

Um campo axissimétrico se decompõe em parte poloidal (nos planos meridionais) e
toroidal (azimutal). A condição $\nabla\cdot\mathbf{B}=0$ é resolvida
automaticamente escrevendo o poloidal em termos de uma \emph{função de fluxo}
$u(\varpi,z)$:
\begin{equation}
\mathbf{B}_{\rm pol} = \frac{\nabla u \times \hat{e}_\phi}{\varpi},
\qquad
\mathbf{B}_{\rm tor} = B_\phi\,\hat{e}_\phi .
\end{equation}
As linhas de campo poloidal são as curvas $u = $ const. Isso é conveniente e
vai reaparecer o tempo todo.

\subsection{Grad--Shafranov}

Impondo que a força de Lorentz seja um gradiente --- condição para
\eqref{eq:bernoulli} existir --- chega-se à equação de Grad--Shafranov:
\begin{equation}
\boxed{\;\Delta^{*}u = -4\pi\varpi^2\rho\,\mu(u) - \beta(u)\beta'(u)\;}
\label{eq:gs}
\end{equation}
onde $\Delta^{*}$ é o operador de Grad--Shafranov e
\begin{equation}
\Delta^{*} \equiv \frac{\partial^2}{\partial\varpi^2}
- \frac{1}{\varpi}\frac{\partial}{\partial\varpi}
+ \frac{\partial^2}{\partial z^2}.
\end{equation}

\begin{alerta}
\textbf{As duas funções arbitrárias.} $\mu(u)$ e $\beta(u)$ \emph{não} são
determinadas pela física --- são escolhas. $\beta(u)$ fixa o campo toroidal
($\varpi B_\phi = \beta(u)$) e $\mu(u)$ a distribuição de corrente poloidal.
Escolhas diferentes dão estrelas diferentes, todas equilíbrios válidos. Boa
parte da literatura --- e boa parte do nosso trabalho --- é sobre \emph{quais
escolhas são fisicamente sensatas}.
\end{alerta}

A escolha mais simples, $\mu = $ const., produz um dipolo de vácuo no exterior.
Fujisawa, Yoshida \& Eriguchi (2012) mostraram que uma lei de potência
negativa,
\begin{equation}
\mu(u) \propto (u+\epsilon)^{m}, \qquad m < 0,
\end{equation}
concentra a corrente perto do eixo e \emph{amplia a região de linhas fechadas},
chegando a campos interiores duas ordens acima do de superfície. Como a equação
passa a ser não linear ($\mu$ depende de $u$), resolve-se iterando.

\ex{Por que $\mu = $ const.\ dá dipolo de vácuo? Dica: pense em que momento de
multipolo uma distribuição de corrente de um só sinal carrega.}

\subsection{Confinamento do toroidal}

Há uma restrição física importante. O campo toroidal precisa \emph{anular-se em
toda linha de campo que chega ao vácuo} --- caso contrário ele vazaria para
fora da estrela, onde não há corrente para sustentá-lo. Isso confina $B_\phi$ à
\textbf{região de linhas fechadas}.

Essa restrição parece técnica e é decisiva. No nosso trabalho ela ocupa
$38$--$45\%$ do volume estelar, e espremer nela a energia necessária para
sustentar massa é exatamente o que quebra o critério (a) da §2.3.

%============================================================
\part*{Parte III --- Estabilidade}
\addcontentsline{toc}{part}{Parte III --- Estabilidade}

\section{Por que campo puro é instável}

Esta é a parte da teoria que mais gente ignora e que mais decide resultados.

\subsection{A instabilidade de Tayler}

Um campo \textbf{puramente toroidal} é instável ao \emph{kink} de Tayler, um
modo $m=1$ análogo à instabilidade de kink em plasmas de laboratório. A força
motriz é a curvatura das linhas de campo: um tubo de fluxo toroidal pode baixar
sua energia magnética deslocando-se lateralmente, e nada o impede.

A taxa de crescimento é da ordem da frequência de Alfvén,
\begin{equation}
\sigma \sim \Omega_A = \frac{v_A}{R},
\qquad
v_A = \frac{B}{\sqrt{4\pi\rho}} ,
\end{equation}
ou seja, o campo é destruído em \emph{poucos tempos de Alfvén}. Para os nossos
parâmetros isso é de ordem de segundos.

\subsection{E poloidal puro também}

Simetricamente, um campo \textbf{puramente poloidal} é instável perto da
superfície neutra, onde as linhas se fecham (Flowers \& Ruderman 1977; a versão
moderna está em Braithwaite \& Spruit).

\subsection{O toro torcido}

Se as duas geometrias puras são instáveis, o que sobra? A resposta, antecipada
por Prendergast (1956) e estabelecida numericamente por Braithwaite \& Spruit
(2006), é a configuração de \textbf{toro torcido}: poloidal e toroidal com
energias \emph{comparáveis}, cada um estabilizando a instabilidade do outro.

\begin{caixa}
\textbf{O critério operacional que usamos.}
\begin{equation}
0{,}1 < \frac{E_{\rm tor}}{E_{\rm pol}} < 10 .
\end{equation}
Os dois ramos puros ficam de fora, como devem.
\end{caixa}

\subsection{A ressalva que você precisa carregar desde o primeiro dia}

\begin{alerta}
O resultado de Braithwaite \& Spruit vale para estrelas \textbf{estavelmente
estratificadas}. A nossa EOS é \textbf{barotrópica} ($P = P(\rho)$, sem
gradiente de entropia ou composição), e há literatura argumentando que
estrelas barotrópicas \emph{não admitem equilíbrio magnético estável nenhum}.

Ou seja: chamar a faixa de energias comparáveis de ``ramo estável'' é uma
conveniência de linguagem no nosso contexto, não um teorema. Essa distinção
foi levantada por uma colaboradora do projeto e está hoje entre as limitações
declaradas do artigo em preparação. Se você quiser um problema aberto com que
contribuir, é um bom candidato.
\end{alerta}

\ex{Por que a estratificação estável estabiliza campos fracos? Pense em uma
parcela de fluido deslocada radialmente numa estrela com gradiente de entropia
positivo e no trabalho que a flutuabilidade faz contra o deslocamento.}

\section{O parâmetro beta}

O critério (a) da §2.3 se escreve com o \emph{beta do plasma}:
\begin{equation}
\boxed{\;\beta \equiv \frac{P_{\rm gas}}{P_{\rm mag}}
= \frac{8\pi P_{\rm gas}}{B^2}\;}
\end{equation}

Onde $\beta < 1$, a pressão magnética domina. Como o equilíbrio foi construído
com uma EOS que não sabe do campo, uma região com $\beta < 1$ é uma região onde
a construção contradiz a si mesma. Por isso exigimos
\begin{equation}
\beta_{\rm min} \ge 1
\end{equation}
sobre o interior estelar ($\rho > 10^6$ g/cm$^3$, para não medir a superfície).

\begin{alerta}
$\beta_{\rm min}\ge1$ é \textbf{consistência da construção, não prova de
estabilidade}. Quem decide estabilidade é evolução dinâmica. Não confunda os
dois --- é um erro fácil e caro.
\end{alerta}

\ex{Para $\rho = 10^9$ g/cm$^3$ e $\mu_e=2$, calcule $P_{\rm gas}$ pela
Eq.~\eqref{eq:eos} e determine qual campo $B$ dá $\beta=1$. Compare com o campo
crítico de Landau $B_c = 4{,}414\times10^{13}$\,G.}

%============================================================
\part*{Parte IV --- Evolução}
\addcontentsline{toc}{part}{Parte IV --- Evolução}

\section{MHD ideal}

Construir o equilíbrio é metade do trabalho. A outra metade é \emph{deixá-lo
evoluir} e ver se sobrevive. As equações são as da magnetoidrodinâmica ideal:
\begin{align}
\frac{\partial\rho}{\partial t} + \nabla\cdot(\rho\mathbf{v}) &= 0, \\
\rho\frac{D\mathbf{v}}{Dt} &= -\nabla P - \rho\nabla\Phi
  + \frac{(\nabla\times\mathbf{B})\times\mathbf{B}}{4\pi},\\
\frac{\partial\mathbf{B}}{\partial t} &= \nabla\times(\mathbf{v}\times\mathbf{B}),
\qquad \nabla\cdot\mathbf{B}=0 .
\end{align}
``Ideal'' significa resistividade nula: o campo é \emph{congelado} no fluido.

\subsection{Escalas de tempo}

Três tempos organizam a leitura de qualquer simulação nossa:
\begin{align}
t_{\rm din} &\sim 1/\sqrt{G\bar\rho} &&\text{(dinâmico, $\sim 0{,}1$--$1$\,s)}\\
t_A &= R/v_A &&\text{(Alfvén, $\sim 2$\,s no nosso modelo)}\\
t_{\rm rot} &= 2\pi/\Omega &&\text{(rotação, $\sim 4$\,s na superfície)}
\end{align}

\ex{Calcule $t_A$ para $B = 3{,}2\times10^{13}$\,G, $\rho = 10^9$\,g/cm$^3$ e
$R = 4\times10^8$\,cm. Se a instabilidade de Tayler cresce em poucos $t_A$,
quanto tempo de simulação você precisa para vê-la?}

\subsection{Transporte restrito e $\nabla\cdot\mathbf{B}=0$}

Numericamente, manter $\nabla\cdot\mathbf{B}=0$ não é automático: erros de
truncamento geram monopolos magnéticos que produzem forças espúrias. A solução
usada no Castro é o \emph{constrained transport}, que armazena $\mathbf{B}$ nas
faces das células e atualiza via forças eletromotrizes nas arestas, preservando
a divergência nula à precisão de máquina por construção.

Por isso exportamos os modelos como \emph{potencial vetor} $\mathbf{A}$ e não
como $\mathbf{B}$: assim $\mathbf{B} = \nabla\times\mathbf{A}$ tem divergência
nula por identidade.

\section{A MRI, e por que não a resolvemos}

A \textbf{instabilidade magnetorrotacional} (MRI) aparece sempre que há rotação
diferencial e campo magnético, e é o mecanismo que apaga rotação diferencial em
discos e em remanescentes de fusão. O comprimento de onda mais instável é
\begin{equation}
\lambda_{\rm MRI} = \frac{2\pi v_{A,z}}{\Omega},
\end{equation}
construído com a componente \emph{vertical} do campo.

\begin{alerta}
\textbf{Esta é a maior limitação do projeto e você precisa sabê-la de cor.}
No nosso modelo de produção $\lambda_{\rm MRI}$ vale cerca de $0{,}004$ de uma
célula da malha. Ou seja: a MRI não é resolvida, e o grupo mostrou que ela não
é resolvível a nenhum custo acessível --- as exigências de que o ambiente
suporte o campo e de que $\lambda_{\rm MRI}$ cubra seis células se fecham em
direções opostas.

Consequência honesta: quando dizemos ``a rotação diferencial sobrevive'', isso
é indistinguível de ``nossa simulação não tem como destruí-la''. Está escrito
assim no artigo, na primeira página, e não nas ressalvas do fim.
\end{alerta}

\section{Resistividade numérica}

Nas equações não há resistividade, mas no \emph{código} há: erros de truncamento
dissipam campo em escalas próximas à da célula. O número de Reynolds magnético
efetivo, $R_m = vL/\eta_{\rm num}$, é grande mas finito e não controlado.

A separação que salva metade dos resultados:
\begin{itemize}
\item A instabilidade de Tayler é \textbf{ideal}. O crescimento, o caráter
      $m=1$ e o transporte de momento angular enquanto o campo existe são
      robustos.
\item O \textbf{decaimento}, o resíduo final e qualquer recrescimento são
      dissipativos e devem ser reportados como limitados por resolução.
\end{itemize}

\ex{Nos nossos runs o mínimo da energia magnética difere por um fator $25$
entre malhas $192^3$ e $256^3$, enquanto a fase de destruição concorda. Explique
por que isso é exatamente o que a separação acima prevê.}

%============================================================
\part*{Parte V --- Este projeto}
\addcontentsline{toc}{part}{Parte V --- Este projeto}

\section{O que medimos}

Uma configuração só conta como resposta se passa \textbf{quatro critérios
simultâneos}:

\begin{center}
\begin{tabular}{ll}
\hline
$M > \Mch = 1{,}44\,\Msun$ & é super-Chandrasekhar de fato \\
$\beta_{\rm min} \ge 1$ & a estrela segura o próprio campo \\
$0{,}1 < E_{\rm tor}/E_{\rm pol} < 10$ & geometria do ramo estável \\
$T/|W| < 0{,}14$, sem perda de massa & rotação que ela consegue manter \\
\hline
\end{tabular}
\end{center}

O quarto é o que torna a resposta honesta: uma estrela sustentada por uma
rotação que dispara o modo barra não é resultado.

\section{Resultados principais, em uma página}

\begin{enumerate}
\item \textbf{Existe um teto.} Uma anã branca segura no máximo
      $E_{\rm tor}/|W| \simeq 0{,}0101$ em campo toroidal confinado, quase
      independente da densidade central. A literatura de equilíbrio precisa de
      $0{,}203$ --- vinte vezes mais.
\item \textbf{O campo compra pouca massa.} Cerca de $+0{,}017\,\Msun$. A
      fronteira de massa converge para $\Mch$ por baixo sem cruzá-la.
\item \textbf{Com rotação, a configuração existe.} Há configurações acima de
      $\Mch$ com energias comparáveis, até $2{,}13\,\Msun$ --- mas a massa é
      \emph{rotacional} e o campo continua em $\sim1\%$ do virial.
\item \textbf{A configuração da literatura perde amplitude.} Uma evolução 3D
      remove $98\%$ da energia magnética em poucos $t_A$, deixando o campo
      ordenado e axissimétrico, porém seis vezes mais fraco.
\end{enumerate}

\section{As ferramentas, e por onde começar}

\begin{center}
\begin{tabular}{ll}
\hline
\code{scf/} & construção de equilíbrios, Python \\
\code{investigations/} & varreduras e análise \\
\code{castro\_problems/} & configuração das simulações 3D \\
\code{tools/} & diagnósticos em C++ sobre plotfiles \\
\code{cluster/} & scripts de submissão (PBS e SLURM) \\
\code{docs/teoria.md} & a referência técnica completa \\
\code{DIARIO.md} & o registro de tudo, com os erros \\
\hline
\end{tabular}
\end{center}

\begin{caixa}
\textbf{Roteiro sugerido para as primeiras duas semanas.}
\begin{enumerate}
\item Resolva os exercícios dos Capítulos 1 e 2 no papel.
\item Leia \code{scf/eos.py} inteiro e confira as fórmulas contra o
      Capítulo 1.
\item Rode o SCF sem campo, reproduza $M(\rho_c)$ e compare com $\Mch$.
\item Ligue a rotação e reproduza $T/|W| = 0{,}099$ do modelo de produção.
\item Leia as §7 e §8 do \code{DIARIO.md} --- são os erros cometidos, e
      valem mais que os acertos.
\end{enumerate}
\end{caixa}

\begin{alerta}
\textbf{Uma regra do grupo, aprendida caro.} \emph{Verificar no destino, nunca
na origem.} ``O arquivo existe aqui'' não é evidência de que existe lá, e numa
cadeia estação $\to$ cluster $\to$ diretório de execução há vários lugares onde
ele pode parar. Um portão custa segundos; uma janela de fila custa horas.
\end{alerta}

%============================================================
\section*{Bibliografia comentada}
\addcontentsline{toc}{section}{Bibliografia comentada}

\paragraph{Fundamentos.}
\begin{itemize}[leftmargin=1.2em]
\item \textbf{Chandrasekhar (1935)}, MNRAS 95, 207 --- o artigo original do
      limite. Curto e legível.
\item \textbf{Shapiro \& Teukolsky}, \emph{Black Holes, White Dwarfs and
      Neutron Stars} --- capítulos 2--4. O livro-texto padrão para tudo na
      Parte I.
\item \textbf{Hachisu (1986)}, ApJS 61, 479 --- o método SCF que usamos.
\end{itemize}

\paragraph{Campos magnéticos em estrelas.}
\begin{itemize}[leftmargin=1.2em]
\item \textbf{Prendergast (1956)}, ApJ 123, 498 --- a primeira configuração de
      toro torcido.
\item \textbf{Braithwaite \& Spruit (2006)}, A\&A (astro-ph/0510316) --- o
      trabalho que estabelece que só a geometria mista é estável. Leitura
      obrigatória.
\item \textbf{Fujisawa, Yoshida \& Eriguchi (2012)}, MNRAS
      (arXiv:1204.5830) --- a prescrição de confinamento de corrente.
\item \textbf{Ciolfi \& Rezzolla (2013)}, arXiv:1306.2803 --- as duas funções
      arbitrárias tratadas com cuidado.
\end{itemize}

\paragraph{Anãs brancas super-Chandrasekhar --- a literatura-alvo.}
\begin{itemize}[leftmargin=1.2em]
\item \textbf{Subramanian \& Mukhopadhyay (2015)}, MNRAS 454, 752 --- a
      família de equilíbrios de onde vem a nossa configuração.
\item \textbf{Gupta, Mukhopadhyay \& Tout (2020)}, MNRAS 496, 894.
\item \textbf{Deb, Mukhopadhyay \& Weber (2022)}, ApJ 926, 1.
\item \textbf{Zuraiq et al.\ (2026)}, ApJL, arXiv:2609.22437 --- o canal de
      formação, via evolução estelar 1D.
\end{itemize}

\paragraph{O que o nosso trabalho acrescenta, e onde ele é frágil.}
\begin{itemize}[leftmargin=1.2em]
\item \code{docs/teoria.md} --- referência técnica completa, com ponteiros
      linha a linha para o código.
\item \code{DIARIO.md} --- o registro cronológico. Leia os erros.
\item \code{paper/ms.tex} --- o artigo em preparação.
\end{itemize}

\vfill
\noindent\rule{\textwidth}{0.4pt}\\
{\small Dúvidas, correções e exercícios resolvidos: converse com o
orientador. Esta apostila é um documento vivo --- se algo aqui estiver errado
ou obscuro, isso é um \emph{bug} e deve ser reportado como tal.}

\end{document}
